\emph{Setting and literature.} The BF reading of MacDowell-Mansouri gravity is due to Freidel and Starodubtsev (hep-th/0501191; cited from memory, as is their quoted coupling $\alpha = G\Lambda/3$ and their normalisation of the BF term); nothing below uses their paper as a premise: the functional $S_{a,b}$ is defined in the statement, the functional it is compared with is that of Fact \claimref{dark-energy.imported-macdowell-mansouri-theorem-stelle-west} in its frozen-point form, and every step is computed on the series' objects. Work in the flat chart of the conformal model, $\Cl(2,4)$ with generators $e_{0}$ ($+1$), $e_{1}, e_{2}, e_{3}$ ($-1$), $e_{+}$ ($+1$), $e_{-}$ ($-1$), $\nO = (e_{+} + e_{-})/2$, $\ninf = e_{+} - e_{-}$, $I = \ninf + (\Lambda/6)\nO$, $X = \nO$, $V = X - (3/\Lambda)I$ (Lemma \claimref{dark-energy.stelle-west-vector-frozen-point}: $\ip{V}{I} = 0$, $Q(V) = -3/\Lambda$, $Q(I)Q(V) = -1$). The fibre volume form is $\epsilon(a, b, P, Q) = \langle a\,b\,P\,Q\,I_{6}^{-1}\rangle_{0}$, as in Theorem \claimref{dark-energy.frozen-point-makes-stelle-west-vector}. A bivector-valued two-form is given by its values $B_{P}$ on the six coordinate planes $P = (m, n)$, $m < n$, and the four-form $B \wedge C$ is $\sum_{P, Q}\omega(P, Q)\,B_{P}C_{Q}\,d^{4}x$ with $\omega(P, Q)$ the sign of the permutation $P Q$ of $(0,1,2,3)$, zero when $P$ and $Q$ share an index; $\omega$ is symmetric and nondegenerate on the six planes.

\emph{(i)} Since $V - X$ is a multiple of $I$, $I \wedge V = I \wedge X = (\ninf + \tfrac{\Lambda}{6}\nO)\wedge\nO = \ninf\wedge\nO$ for every $\Lambda \neq 0$, and $(\ninf \wedge \nO)^{2} = \ip{\ninf}{\nO}^{2} - Q(\ninf)Q(\nO) = 1$. So $I \wedge V = e_{+}e_{-}$, $\iota = (I\wedge V)I_{6} = e_{0}e_{1}e_{2}e_{3}$ and $\iota^{2} = -1$ (computed at $\Lambda = 3/4, -3/2, 12$). For $I \perp V$ and bivectors $P, Q$, $\epsilon(I, V, P, Q) = \langle (I\wedge V)\,P\,Q\,I_{6}^{-1}\rangle_{0}$, and on the fifteen basis bivectors $\epsilon(I, V, P, Q) = -\langle P\,\iota\,Q\rangle_{0}$ (computed, all $225$ pairs); the form is symmetric (reversion), invariant under the six generators of the stabiliser of $(I, V)$ (computed as the kernel of $B \mapsto (B\times I, B\times V)$, dimension $6$ as in Theorem \claimref{dark-energy.frozen-point-makes-stelle-west-vector}(ii)), and odd under the internal reflection in $e_{1}$, which reverses $\iota$ (internal only: as spacetime densities the $\epsilon$ terms are even under ordinary parity). Orientation: in this chart $\epsilon(I, V, e\wedge e, e\wedge e) = +6\det(\hb^{-1})$ (part (v)); with the opposite orientation $\nO = (e_{+} - e_{-})/2$, $\ninf = e_{+} + e_{-}$ one has $\ninf\wedge\nO = -e_{+}e_{-}$, so every value of $\epsilon(I, V, \cdot, \cdot)$ changes sign, and a sum over all $24$ orderings of the four form indices is four times the plane-sum coefficient used here. The six slots of $\epsilon$ hold one vector and two bivectors only if a second vector fills the sixth (the slot count of Theorem \claimref{dark-energy.frozen-point-makes-stelle-west-vector}(iii)), so the frozen point alone cannot carry this term.

\emph{(ii)} Fix a point and write $L(B) = \sum\omega(P, Q)\big[\langle B_{P}F_{Q}\rangle_{0} - \tfrac14\langle B_{P}\,M B_{Q}\rangle_{0}\big]$ with $M = a\iota + b$ acting on the Lorentz block (and zero elsewhere). Because $\iota$ commutes with the bivectors of $\{I, V\}^{\perp}$ and $\langle P\iota Q\rangle_{0} = \langle Q\iota P\rangle_{0}$ (reversion), the quadratic part is symmetric and its gradient in the direction $\delta B_{P_{0}} = d$ is $\sum_{Q}\omega(P_{0}, Q)\big[\langle d\,F_{Q}\rangle_{0} - \tfrac12\langle d\,M B_{Q}\rangle_{0}\big]$. For $d$ off the Lorentz block (in $V \wedge \{I, V\}^{\perp}$), $\langle d\,M B_{Q}\rangle_{0} = 0$ (the product keeps a factor $V$), so the gradient is $\sum_{Q}\omega(P_{0}, Q)\langle d\,F_{Q}\rangle_{0}$, independent of $B$; since $\omega$ and the restriction of $\langle\cdot\,\cdot\rangle_{0}$ to $V \wedge \{I, V\}^{\perp}$ are nondegenerate, these vanish for all such $d$ and $P_{0}$ exactly when the translational block of $F$ is zero, which for the reduced connection $\Omega(a) - \hb^{-1}(a)\wedge V/\ell^{2}$ is the torsion $De$ (Theorem \claimref{dark-energy.frozen-point-makes-stelle-west-vector}(ii)); the translational components of $B$ then do not appear in $L$ and are multipliers. On the Lorentz block, $M$ is invertible with $M^{-1} = (b - a\iota)/(a^{2} + b^{2})$ (because $\iota^{2} = -1$), and the gradient vanishes for $F_{Q} = \tfrac12 M B_{Q}$, that is $B_{L} = 2M^{-1}F_{L}$. Uniqueness: the quadratic form $(P, l), (Q, m) \mapsto \omega(P, Q)\langle l\,\iota\,m\rangle_{0}$ on the $36$ Lorentz-block components has rank $36$ (computed), and for $b \neq 0$ the same holds since $M$ is invertible. The value is $L = 2\langle M^{-1}F\wedge F\rangle_{0} - \langle M^{-1}F\wedge F\rangle_{0} = \langle M^{-1}F \wedge F\rangle_{0} = (a\,\epsilon(I, V, F, F) + b\,\langle F\wedge F\rangle_{0})/(a^{2} + b^{2})$, using $-\langle F\iota F\rangle_{0} = \epsilon(I, V, F, F)$ from (i). This is checked exactly on a fixed integer table of Lorentz-block field strengths for four pairs $(a, b)$, including $a = 0$ and $a < 0$, with the gradient zero in all $90$ directions, and for a field strength with a translational component the gradient in that direction is the same nonzero number for two different $B$. With coefficient $\lambda$ on the BF term the stationary point is $\lambda$ times the one above and the value is $\lambda^{2}$ times it (checked at $\lambda = 3$): BF coefficient $\lambda$ with breaking $(a, b)$ gives the same effective functional as BF coefficient $1$ with $(a, b)/\lambda^{2}$.

For the part $b$, even under the internal reflection, on the Lorentz block $F = R + \tfrac{\Lambda}{3}\,e\wedge e$, so $\langle F \wedge F\rangle_{0} = \langle R\wedge R\rangle_{0} + \tfrac{2\Lambda}{3}\langle R \wedge e\wedge e\rangle_{0} + \tfrac{\Lambda^{2}}{9}\langle e\wedge e\wedge e\wedge e\rangle_{0}$. The last is $\sum\omega(P,Q)\langle (e_{m}\wedge e_{n})(e_{p}\wedge e_{q})\rangle_{0}$ over complementary planes, zero because complementary coordinate bivectors have no common vector (computed zero). The middle term, in components $\epsilon^{mnpq}R_{mnpq}$ up to a constant, vanishes by the first Bianchi identity of a torsion-free field (computed zero on the de Sitter and Bianchi I samples). The first is the Pontryagin density of the Lorentz connection, closed by the Chern-Weil theorem (standard), so it does not contribute to the field equations; and at $a = 0$ the whole value is $\langle F\wedge F\rangle_{0}/b$, topological on the constraint surface.

\emph{(iii)} Vary the Lorentz connection by $\delta\Omega$ in $\int\epsilon(I, V, F, F)$: $\delta F_{L} = D\delta\Omega$, and after an integration by parts the variation is $-2\int\epsilon(I, V, \delta\Omega, DF_{L})$ with $DF_{L} = DR + \tfrac{\Lambda}{3}D(e\wedge e) = \tfrac{2\Lambda}{3}\,T\wedge e$ by the Bianchi identity $DR = 0$, $T = De$. The connection equation is therefore the linear map $T \mapsto \big(b, q \mapsto \sum\epsilon^{qmnp}\epsilon(I, V, b, T_{mn}\wedge e_{p})\big)$ from the $24$ components of $T$ to the $24$ pairs (Lorentz bivector $b$, direction $q$); its rank is $24$ (computed), so in vacuum it is equivalent to $T = 0$. For $S_{\alpha, 0}$ with $B_{L}$ eliminated the field equations are those of $\alpha^{-1}\int\epsilon(I, V, F, F) + \int\langle B_{V}\wedge F\rangle_{0}$, $B_{V}$ the translational multiplier: $T = 0$; the connection equation $-2\alpha^{-1}\epsilon(I, V, \cdot, DF_{L}) + (\text{Lorentz part of } [\,\hb^{-1}\wedge V/\ell^{2}, B_{V}])= 0$; and the frame equation of the functional plus $D B_{V}$. With $T = 0$ the first term of the connection equation vanishes, and the second is the map $\beta \mapsto \sum\epsilon^{qmnp}\beta^{a}_{mn}\,e_{p}\wedge e_{a}$ (from $[e_{p}\wedge V, e_{a}\wedge V] = -Q(V)\,e_{p}\wedge e_{a}$), which is the map just computed composed with the invertible $\iota$; so $B_{V} = 0$ and the frame equation is that of the functional. Conversely every vacuum solution of the functional has $T = 0$ and gives a stationary point with $B_{V} = 0$, $B_{L} = -(2/\alpha)\iota F_{L}$. With a matter functional coupled to the connection, its spin current enters the connection equation; $T = 0$ still holds (it is the $B_{V}$ equation, independent of matter), and the rank-$24$ map solves $B_{V}$ for the spin current, which then enters the frame equation through $DB_{V}$. This is the torsion-free coupling with the symmetric (Hilbert, Belinfante) energy-momentum tensor, whereas the stored functional with the Einstein-Cartan coupling of Fact \claimref{synthesis.imported-einstein-cartan-matter-coupling} has torsion fixed algebraically by the spin; the two differ by the spin-spin contact interaction of order $G$ times the square of the spin density (Hehl, von der Heyde, Kerlick and Nester 1976, cited from memory).

\emph{(iii$'$)} The form $\langle B_{V}B'_{V}\rangle_{0}$ on $V \wedge \{I, V\}^{\perp}$ is, up to sign, $Q(V)$ times the metric of $\{I, V\}^{\perp}$, so it is nondegenerate and invariant under the stabiliser of $(I, V)$ (checked on its six generators). In $S_{a,b,c}$ the gradient in a translational direction $d$ at plane $P_{0}$ is $\sum_{Q}\omega(P_{0}, Q)\big[\langle d\,F_{Q}\rangle_{0} - \tfrac{c}{2}\langle d\,B_{V,Q}\rangle_{0}\big]$, zero for $B_{V} = 2F_{V}/c$, and the $B_{V}$ part of the value is $\tfrac{2}{c}\langle F_{V}\wedge F_{V}\rangle_{0} - \tfrac{c}{4}\cdot\tfrac{4}{c^{2}}\langle F_{V}\wedge F_{V}\rangle_{0} = \tfrac1c\langle F_{V}\wedge F_{V}\rangle_{0}$ (checked exactly at $c = 3$ and $c = -1/2$ on a field strength with a translational block, where $\sum\omega\langle F_{V}F_{V}\rangle_{0} \neq 0$). $F_{V}$ is the torsion block, so this is a torsion-squared term; it replaces the constraint $T = 0$, and as $c \to \infty$ the effective functional tends to $(a\int\epsilon(I, V, F, F) + b\int\langle F_{L}\wedge F_{L}\rangle_{0})/(a^{2}+b^{2})$ on the unconstrained field space of connections preserving $I$, which at $b = 0$ is the stored functional, whose connection equation with spinning matter gives the Einstein-Cartan torsion. The even term of Freidel and Starodubtsev (from memory) may be the Killing form on the whole de Sitter algebra, which contains this translational piece; if so, their action and $S_{a,b}$ differ exactly in the matter sector.

\emph{(iv)} At $a = b = 0$, $S$ is $\int\langle B \wedge F\rangle_{0}$, whose $B$-equation is $F = 0$ and $A$-equation $DB = 0$. By Theorem \claimref{synthesis.gtg-tractor-gauge-field-reproduces-tractor-connection} the tractor gauge field on the Einstein region has algebraic field strength $-W$, which annihilates $I$ and $V$ (Theorem \claimref{dark-energy.frozen-point-makes-stelle-west-vector}(i), (ii)), so its translational block vanishes, so $F = 0$ is $W = 0$ with the Einstein condition: the locally de Sitter or anti-de Sitter fields, where tractor transport is path-independent. Checked: on the de Sitter field ($H = 1/2$, $\Lambda = 3/4$) the Lorentz block $-\mathcal{R}(e_{m}\wedge e_{n}) + \tfrac{\Lambda}{3}e_{m}\wedge e_{n}$ is zero on all six planes, so $B = 0$. With $a = \alpha$, $b = 0$ and a vacuum Einstein field, (ii) gives $B_{L} = -(2/\alpha)\iota F_{L}$ and $F_{L} = -W$, so $B_{L} = (2/\alpha)\iota W$; checked with $\mathcal{R} = \tfrac{\Lambda}{3}\,\mathrm{id} + $ the Weyl operator of the Kasner vacuum $(\tfrac23, \tfrac23, -\tfrac13)$ at $\tau = 8$ (the construction of the stored frozen-point check; the Kasner Ricci map is checked zero): $B_{L}$ is nonzero on all six planes and commutes with $I$ and $V$.

\emph{(v)} On a synchronous diagonal field, with $R_{P} = -\mathcal{R}(e_{m}\wedge e_{n})$ and $(e\wedge e)_{P} = e_{m}\wedge e_{n}$, $e_{m} = \hb^{-1}(e_{m})$, the sums $\sum\omega\,\epsilon(I, V, R_{P}, (e\wedge e)_{Q})$ and $\sum\omega\,\epsilon(I, V, (e\wedge e)_{P}, (e\wedge e)_{Q})$ equal $-\tfrac12\mathcal{R}\det(\hb^{-1})$ and $6\det(\hb^{-1})$ on the de Sitter field and on a generic Bianchi I field (data $(2, 3, -1)$, $(\tfrac13, 2, 5)$, $(3, -\tfrac12, 7)$ for $(a_{i}, \dot a_{i}, \ddot a_{i})$), at three values of $\Lambda$ in $I, V$ (computed; $\mathcal{R} = 4\Lambda$ on de Sitter fixes the store's sign). The first sum is a Lorentz-invariant linear function of the field strength (the curvature map) at the point; the only such function is a multiple of $\mathcal{R}$ (the Weyl and trace-free Ricci parts are nontrivial irreducible representations), so one sample with $\mathcal{R} \neq 0$ fixes the multiple and the second sample corroborates it. Expanding $F = R + \tfrac{\Lambda}{3}e\wedge e$ as in Theorem \claimref{dark-energy.frozen-point-makes-stelle-west-vector}(iii), $\tfrac{1}{Q(I)}\epsilon(I, V, F, F) = \tfrac{3}{\Lambda}\epsilon(R, R) + 2\epsilon(R, ee) + \tfrac{\Lambda}{3}\epsilon(ee, ee) = \tfrac{3}{\Lambda}\epsilon(I, V, R, R) - (\mathcal{R} - 2\Lambda)\det(\hb^{-1})$ (checked as an identity on both samples at three values of $\Lambda$). The Einstein-Hilbert and cosmological parts thus enter as $-(\mathcal{R} - 2\Lambda)\sqrt{|g|}$ with coefficient of magnitude one. In Theorem \claimref{synthesis.macdowell-mansouri-principle-fixes-the-trace-free-part}(iv) the functional is normalised so that its frame equation reads $\mathcal{G}(a) + \Lambda a + 8\pi G\,T(a) = 0$, with $G$ the normalisation constant, which is the Einstein-Hilbert normalisation $(16\pi G)^{-1}\int(\mathcal{R} - 2\Lambda)\sqrt{|g|}$ up to an overall sign convention; so the stored functional is $k\,Q(I)^{-1}\int\epsilon(I, V, F, F)$ with $|k| = 1/(16\pi G)$. Comparing with (ii), $|a|/(a^{2} + b^{2}) = 1/(16\pi G\,|Q(I)|)$. With $N = M_{P}/H$, $G = M_{P}^{-2}$ and $Q(I) = H^{2}$ (Fact \claimref{standard-model.imported-frozen-point-dictionary}; $N^{2} = M_{P}^{2}/Q(I)$ in Theorem \claimref{synthesis.de-sitter-entropy-pi-n-squared}), $1/(G\,Q(I)) = N^{2}$, so $|a| N^{2} = 16\pi a^{2}/(a^{2} + b^{2})$, and at $b = 0$, $|\alpha| = 16\pi G\,Q(I) = 16\pi/N^{2}$, about $1.6 \times 10^{-121}$ for $N \approx 1.0 \times 10^{61}$. By the rescaling of (ii) the same effective functional is reached with BF coefficient $\lambda$ and $\alpha/\lambda^{2}$, and the factor $16\pi$ is replaced by $1$ in the normalisation Freidel and Starodubtsev quote (from memory); so $\alpha N^{2}$ is a convention of the BF normalisation and of $b$. This is the pre-registered test of a relation between the coupling and $N$, and it is negative: $\alpha$ is $N^{-2}$ re-expressed, dimensional analysis having already left $\hbar G\Lambda$ as the only dimensionless combination.

\emph{Reading (argued).} That the frozen point alone admits no invariant of this kind, odd under the internal reflection, is the content of the classification of invariant quadratic forms, a separate result, and is not used here. The frozen point's direction reduces the ambient group (Theorem \claimref{dark-energy.frozen-point-makes-stelle-west-vector}(i)) and its norm sets the coupling, while the breaking is on the blade $I \wedge X$. The MacDowell-Mansouri form is rewritten here, not reduced: choosing the breaking $\epsilon(I, V, B, B)$ is the same choice as the integrand $\epsilon(I, V, F, F)$, since elimination inverts the quadratic form, and the field space (connections preserving $I$) is the stored functional's. The action principle itself remains an input, and the matter coupling is decided by the support of the breaking ((iii), (iii$'$)), not by the BF form. The Euler term is part of $\epsilon(I, V, F, F) = \epsilon(I, V, R, R) + \tfrac{2\Lambda}{3}\epsilon(I, V, R, e\wedge e) + \tfrac{\Lambda^{2}}{9}\epsilon(I, V, e\wedge e, e\wedge e)$ and comes out with it, with coefficient $1/\alpha = N^{2}/(16\pi)$ in $S_{\alpha, 0}$. The frozen-point action keeps $\Lambda$, hence $\alpha$, fixed. If instead the coupling is promoted globally to a field, as in the Henneaux-Teitelboim form (Fact \claimref{dark-energy.imported-henneaux-teitelboim-promotion}, recorded in the store as the variational form of the Clausius route), the variable conjugate to $\alpha$ is $\partial S/\partial\alpha = \tfrac14\int\epsilon(I, V, B, B)$, the integrated breaking term, which on shell equals $-\alpha^{-2}\int\epsilon(I, V, F, F)$ and in vacuum is the integral of $\epsilon(I, V, W, W)$, the Weyl-squared density of the stored time conjugate to the norm (Theorem \claimref{synthesis.promoted-norm-weyl-clock-flat-weight}); a local promotion would instead impose $\epsilon(I, V, F, F) = 0$ pointwise, which Schwarzschild-de Sitter violates. These identifications, with $\ell$ absorbed into the frame, are argued and not checked. Transport-convenient: the unbroken sector is pure transport, and gravity is its breaking on the plane of the frozen point and the canonical tractor.
